% Test babelbidi-api-layout.def; compile twice with any of the three engines.
\documentclass{article}
\usepackage{amsmath}
\usepackage[main=english]{babelbidi}
\newtheorem{apitest}{Theorem}[subsection]

\title{Layout API}
\author{First author\LTRthanks{LTR title footnote.}%
    \and Second author\RTLthanks{RTL title footnote.}}
\date{}

\begin{document}
\maketitle

\section{Footnotes and overlaps}

Separate mark\footnotemark
\LTRfootnotetext{LTR text with the current footnote number.}
Explicit mark\footnotemark[7]
\RTLfootnotetext[7]{RTL text with an explicit number.\par Second paragraph.}

{\setfootnoteRL
    Ordinary RTL\footnote{Ordinary footnote forced RTL.}
    Explicit LTR\LTRfootnote{Explicit LTR overrides the RTL setting.}
    \setfootnoteLR
    Ordinary LTR\footnote{Ordinary footnote forced LTR.}
    Explicit RTL\RTLfootnote{Explicit RTL overrides the LTR setting.}
    \setfootnoteRL\unsetfootnoteRL
    Unset RTL\footnote{LTR after unsetfootnoteRL.}}

Default restored\footnote{The original footnote behaviour is restored.}

\medskip
\begin{minipage}{.8\linewidth}
    \setfootnoteRL
    Minipage\footnote{RTL minipage footnote.}

\end{minipage}

\LRfootnoterule

\bigskip
\noindent LTR overlaps:\quad
\hboxL{\vrule height 10pt\bidillap{L}\bidirlap{R}}

\medskip
\noindent RTL overlaps:\quad
\hboxR{\vrule height 10pt\bidillap{L}\bidirlap{R}}

\clearpage
\section{Lists and bibliography}

\begin{enumerate}
    \item Default LTR item.
    \begin{RTLitems}
        \item RTL item.
        \begin{itemize}
            \item Nested RTL item.
            \begin{LTRitems}
                \item Nested LTR item.
            \end{LTRitems}
        \end{itemize}
    \end{RTLitems}
    \item LTR restored.
\end{enumerate}



\setRTLbibitems
\begin{thebibliography}{9}
    \bibitem{rtl} RTL entry with an LTR heading.
    \begin{LTRbibitems}
        \bibitem{ltr-group} LTR entry in an RTL bibliography.
    \end{LTRbibitems}
    \bibitem{rtl-restored} RTL restored.
\end{thebibliography}

\begin{RTL}
    \setLTRbibitems
    \begin{thebibliography}{9}
        \bibitem{ltr} LTR entry with an RTL heading.
        \begin{RTLbibitems}
            \bibitem{rtl-group} RTL entry in an LTR bibliography.
        \end{RTLbibitems}
    \end{thebibliography}
\end{RTL}

\setdefaultbibitems
\begin{thebibliography}{9}
    \bibitem{default} Default bibliography direction restored.
\end{thebibliography}

\clearpage
\section{Cases and numbering}

\SepMark{-}
\subsection{Hyphen separator}\label{api:hyphen}
\subsubsection{Nested numbering}
\numberwithin{equation}{subsection}

\begin{equation}\label{api:cases}
    f(x)=\begin{RTLcases}
        $x^2$ & \LR{if $x>0$}\\[3pt]
        $-x$  & \LR{otherwise}
    \end{RTLcases}
\end{equation}

Text mode: \begin{RTLcases}
    \LR{first} & \LR{condition A}\\
    \LR{second longer} & \LR{B}
\end{RTLcases}.

\begin{RTL}
    \begin{RTLcases}
        $x+1$ & \LR{condition}\\
        $0$   & \LR{otherwise}
    \end{RTLcases}
\end{RTL}

\begin{apitest}
    The parent number uses the selected separator.
\end{apitest}

References: subsection~\ref{api:hyphen}, equation~\eqref{api:cases}.

{\SepMark{/}}
\subsection{Slash separator persists outside its group}
\begin{equation}a+b=c\end{equation}

\SepMark{.}
\subsection{Default separator restored}
\begin{equation}a-b=c\end{equation}

\clearpage
\section{Margin notes}

Default placement.\marginpar{\LR{Default}}
\par\bigskip

\setRTLmarginpar
RTL policy.\marginpar{\LR{RTL policy}}
\par\bigskip
Explicit LTR text, with the same placement.\LTRmarginpar{\LR{LTR text}}
\par\bigskip

\setLTRmarginpar
LTR policy.\marginpar{\LR{LTR policy}}
\par\bigskip
Explicit RTL text, with the same placement.\RTLmarginpar{\LR{RTL text}}
\par\bigskip

\setdefaultmarginpar
Optional text.\LTRmarginpar[\LR{Left text}]{\LR{Right text}}
\par\bigskip

\reversemarginpar
Reversed margin.\RTLmarginpar[\LR{Left choice}]{\LR{Right choice}}
\par\bigskip
\normalmarginpar

\begin{RTL}
    \LR{Default placement in RTL text.}\marginpar{\LR{RTL default}}
    \par\bigskip
    % Native bidi applies the explicit direction only to the first paragraph;
    % the LuaTeX/pdfTeX implementation keeps it for the whole margin note.
    \LR{Explicit LTR text in RTL context.}\LTRmarginpar{First line.\par Second line.}
    \par\bigskip
\end{RTL}

Default restored.\marginpar{\LR{Restored}}


\clearpage
\RTLdblcol
\twocolumn[\section{Two-column order}]
% Recompute the common rule width after switching to narrow columns.
\setlength{\footnoterulewidth}{.4\columnwidth}
\setlength{\columnseprule}{.4pt}

\noindent RTL: first column (right).\LTRfootnote{Footnote in the first column.}
\par\medskip
The text itself remains left to right.

\newpage
\noindent RTL: second column (left).\LTRfootnote{Footnote in the second column.}

\clearpage
\LTRdblcol
\noindent LTR: first column (left).\LTRfootnote{Footnote in the first column.}

\newpage
\noindent LTR: second column (right).\LTRfootnote{Footnote in the second column.}

\onecolumn
\setlength{\footnoterulewidth}{.4\columnwidth}
\setlength{\columnseprule}{0pt}



\clearpage
\section{Continued footnote rules and diagnostics}
\autofootnoterule
\debugfootnotedirection
\SplitFootnoteRule
\newcommand{\NoteLine}{A line in a long footnote.\par}
\newcommand{\NoteBlock}{%
    \NoteLine\NoteLine\NoteLine\NoteLine\NoteLine\NoteLine
    \NoteLine\NoteLine\NoteLine\NoteLine\NoteLine\NoteLine}
A long note\LTRfootnote{%
    \NoteBlock\NoteBlock\NoteBlock\NoteBlock\NoteBlock\NoteBlock}
\clearpage
RTL note\RTLfootnote{\LR{A right-to-left note.}}
\end{document}
